\section{原创探索与创新点的具体拆解}
\label{sec:innovation}
本节将课程要求中的“原创创新点”落实为可以逐项核查的实验贡献。这里的原创探索，指本次围绕具体问题完成的设计、实现和经验分析，不表示相关思想在数据库领域首次出现。排序改善范围统计的作用已有技术依据，重组成本与查询收益的权衡也已有研究\cite{duckdb,rong}。本文在此基础上，从\textbf{成本口径、窗口位置和误差解释}三个层面细化问题。以下三项均由前文已有记录支撑；本节的补写没有新增性能实验。

\subsection{探索一：把单次加速比较扩展为分条件的累计收益验证}
\label{sec:innovation-cost}
\textbf{已有知识与未回答的问题。}
排序文件查询更快，回答的是“排序文件已经存在时，查询阶段是否省时”。对于手中只有无序文件的使用者，重新读取、排序和写出是一笔先发生的成本。仅报告加速倍数，既不能确定查询多少次后覆盖成本，也不能区分刚好覆盖成本与获得约定幅度的收益。已有成本思想不能直接给出本数据与本负载下的数值边界。

\textbf{本次具体设计。}
本文固定数据内容、行组大小、压缩和主要引擎设置，只保留无序与时间排序两种布局。在主实验同一批次中分别测量一天、七天窗口；每一轮，两布局使用完全相同的查询序列。一次重排成本 $R$ 与整批查询耗时分别实测，再按式\eqref{eq:cost}--\eqref{eq:gain} 合并。每个 $N$ 都重新执行完整查询批次，而不是用单次平均耗时乘 $N$ 充当实测。这使比较对象从“一个查询快多少”变成“这一批任务累计能省多少”。

\textbf{两种边界为何不同。}
为解释判据而作一个简化推导：假设同一稳定负载下，无序和排序的平均每次查询耗时分别为 $u$、$s$，并假设累计查询时间可以近似为 $Nu$、$Ns$。对于预设降幅 $0\le\tau<1$，有
\begin{align}
Nu-(R+Ns)>0 &\iff N(u-s)>R,\label{eq:innovation-zero}\\
\frac{Nu-(R+Ns)}{Nu}\ge\tau
 &\iff N\bigl[(1-\tau)u-s\bigr]\ge R.\label{eq:innovation-target}
\end{align}
当 $R>0$ 时，只有 $u>s$ 才有有限的零收益交点；要达到给定正收益门槛，还需 $(1-\tau)u>s$。条件满足时，连续近似交点分别为 $R/(u-s)$ 和 $R/[(1-\tau)u-s]$。因此，“排序查询较快”并不足以保证任意指定的收益幅度一定可达。取 $\tau=0.10$ 时，达到门槛通常需要比刚覆盖成本更多的查询。

这一推导只是已有成本模型的代数解释，不是新算法，也不是替代实测的拟合结果。本次实际查询序列与耗时存在波动，正文判断仍使用每轮整批记录及离散已测次数，不用上述近似式宣布精确阈值。10\pct{} 也是本次约定，不是行业最优值或统计显著性标准。

\textbf{验证结果与可支持的贡献。}
主实验一天／七天窗口在 \MainOneNetN{}／\MainSevenNetN{} 次时持续满足五轮净节省为正，但对应降幅中位数只有 \MainOneNetPct\pct{}／\MainSevenNetPct\pct{}；首次五轮均达到 10\pct{} 的已测点则是 \MainOneTenN{}／\MainSevenTenN{} 次，见表\ref{tab:thresholds}。贡献在于将不同窗口负载与两类判据放入同一套可核查的累计实验，实际展示了“查询加速、覆盖成本、达到约定收益”三个结论不能混用。

\textbf{改进的对象与适用边界。}
本文改进的是本案例的评估和报告方式，而非 Parquet 读写器：同时报告一次成本、整批查询成本、逐轮净节省和门槛通过轮数。这样的结果更能说明一项优化在给定任务下是否值得，但仍是离线评估。它已测过全量排序成本，不能声称在尚未排序时就能自动作出正确选择，也不能把 \MainOneTenN{} 或 \MainSevenTenN{} 次推广到其他设备和负载。

\subsection{探索二：固定窗口长度，检查行组边界位置的影响}
\label{sec:innovation-boundary}
\textbf{从主实验中发现的具体缺口。}
一天与七天对比同时改变了窗口长度、命中记录数和候选行组数。主实验的一天窗口又全部落在一个候选行组内。因此，它不能回答一个更窄的问题：即使都查连续 24 小时，仅窗口起点相对文件行组边界的位置不同，收益判断是否仍会变化？这项问题是在主实验之后提出，属于有时间顺序记录的补充探索。

\textbf{提出并实施的比较方式。}
选择六个固定行组位置，为每处构造一对 24 小时窗口：一类以行组时间范围中点为中心，另一类以下一行组的起始时间为中心。前者只保留一个候选行组，后者保留两个。选择依据是元数据，不是先看哪个窗口测得更快。每对命中行数差异事先要求不超过 3\pct{}，实际为 0.03\pct--1.49\pct；窗口清单、调度顺序和判断标准在本轮性能计时前冻结。两类比较既有每窗口九轮单次计时，也有每个 $N$ 五轮整批计时，均检查结果一致性，详见第\ref{sec:supp-method}节。

\textbf{为什么这一设计比继续加长窗口更具体。}
固定 24 小时长度，避免把“时间跨度更长”直接当作唯一解释；按相近命中数配对，减轻命中记录规模差异的影响；同一补充批次交替测量，并为每个窗口保留无序基线，避免拿前一天的一组时间与当天另一组时间简单相减。六对窗口也使结论不只来自单个有利例子。不过，这些控制没有让两类日期与数据完全相同，所以它是匹配比较，不是隔离所有变量的因果实验。

\textbf{实际证据及其含义。}
六对窗口中，跨界窗口的排序查询耗时中位数均高于组内窗口；两类汇总值分别为 \SuppCrossQueryMs{} 和 \SuppInsideQueryMs{} 毫秒。整批结果按中位数达到 10\pct{} 的首次已测点为组内 \SuppInsideMedianN{} 次、跨界 \SuppCrossMedianN{} 次；要求五轮全部达标时为 \SuppInsideAllN{} 和 \SuppCrossAllN{} 次。表\ref{tab:paired-saving}保留每对单次节省差值及正差轮数；表\ref{tab:supp-thresholds}同时保留未全部达标的轮数，避免只给出两个整数便暗示一个确定不变的差距。

\textbf{本次探索的价值与限制。}
这一结果把经验建议从“按查询天数判断”细化为“还应检查查询窗口与物理行组的关系”。可采用的评估改进是：测试一天窗口时同时覆盖组内和跨界位置，并报告命中行数、候选组数与配对收益。这里增加的是验证覆盖面，而不是发明新的行组裁剪机制。候选组来自元数据，不等于实测物理读取；跨界与组内使用不同日期，累计比较还沿用历史 $R$。因此只能说在本批样本中观察到关联，不能声称“跨一组必然慢固定比例”，也不能用主、补实验之间的次数差证明边界的独立作用。

\subsection{探索三：拆开总耗时误差，诊断固定成本的遮蔽作用}
\label{sec:innovation-error}
\textbf{问题从何而来。}
原误差复核只评价排序方案总耗时 $R+Q_N$。当一次性排序成本占总耗时大部分时，即使查询阶段估计偏差较大，总耗时相对误差也可能不大。因此，“总误差低于 20\pct{}”不能自动推出查询阶段预测同样准确，更不能推出相对于无序方案的净收益判断准确。

\textbf{本次分析上的改进。}
保留原有总耗时指标 $E_T$，另报告查询阶段指标 $E_Q$；两份预实验参数、两类窗口和全部 19 个 $N$ 都保留，共 76 个组合点，不选择误差最小的参数来源。式\eqref{eq:decompose}进一步将带符号总误差拆成排序成本项和查询项，说明固定成本权重及异号抵消怎样影响总误差。该恒等式是成本加和的代数展开，属于解释工具，不宣称为新定理；$E_Q$ 是事后诊断量，没有补设一个有利的通过阈值。

\textbf{具体例子与可复核结论。}
一天窗口、一次查询时，历史排序成本占排序方案总耗时约 99.16\pct{}；查询阶段误差对总误差的权重很低。第一轮预实验参数在一天窗口、100 次查询时，总耗时误差为 \PilotOneDayOneHundredTotalErr\pct{}，而查询阶段误差为 \PilotOneDayOneHundredQueryErr\pct{}。两个数字来自同一条件，表明只给出前者会掩盖估计中不同组成部分的表现。详细输入、误差方向和其他组合见第\ref{sec:error}节及表\ref{tab:errors}。

\textbf{贡献边界。}
本项提供一个针对本案例的误差诊断与报告改进：同时报告总成本准确度、查询阶段准确度，以及固定成本占比。它没有提高预测器本身的准确率，没有新增性能测量，也没有完成对无序耗时或净收益预测的独立验证。分析发生在正式实验之后，应作为解释性补充，而不能倒写成预先登记且成功通过的新预测方法。

\subsection{三项探索的对应证据与实际贡献归属}
\begin{table}[htbp]
\centering\small
\caption{课程层面原创探索的对象、产出和边界}\label{tab:innovation-map}
\begin{tabularx}{\linewidth}{@{}>{\raggedright\arraybackslash}p{.18\linewidth}YY@{}}
\toprule
探索对象 & 本次完成的具体产出 & 证据与不应扩大的结论\\
\midrule
累计成本与收益判据 & 同序列整批比较；区分净节省和 10\pct{} 收益；按窗口长度报告 & 主实验计时、表\ref{tab:thresholds}；不是新算法或通用回本次数\\
固定长度的边界位置 & 元数据选取六对窗口；命中数匹配；单次与整批两层验证 & 冻结窗口及补充记录、表\ref{tab:single-supp}和\ref{tab:supp-thresholds}；不是唯一因果机制证明\\
总误差的组成 & 保留 76 点；并列总误差、查询误差和成本权重 & 事后误差文件、表\ref{tab:errors}；不是预测精度改进或事前验证\\
\bottomrule
\end{tabularx}
\end{table}

以上三项共同形成“提出窄问题---固定比较条件---实测与校验---解释有效范围”的实验链条。代码、冻结清单、逐轮记录和数据宏让这些判断能够追溯到具体输入和输出；记录完整性是支撑证据的工程工作，本身不作为新的数据库技术。

“本次探索”也不等于每个观点均由学生独立首次提出。选题与方案细化、部分追问和报告审查由两个 AI 对话协助完成；学生通过对话选择采用研究范围与验证要求，Codex 编码、调用工具运行和整理材料。具体建议来源与学生实际选择应以两个对话的完整记录核对，见附录\ref{app:ai}。据此，本文将贡献定位为学生主导任务选择、AI 深度协助的课程实验探索，不将转述的 AI 建议标为个人独立首创。
\FloatBarrier
